\documentclass[a4paper,12pt]{article}
\usepackage[T2A]{fontenc}
\usepackage[utf8]{inputenc}
\usepackage[russian]{babel}
\usepackage{amsmath,amssymb,mathtools}
\usepackage{helvet}
\renewcommand{\familydefault}{\sfdefault}
\usepackage{icomma}
\usepackage[a4paper,left=23mm,right=23mm,top=21mm,bottom=24mm,headheight=14pt]{geometry}
\usepackage{microtype}
\usepackage{tikz}
\usetikzlibrary{arrows.meta,positioning,shapes.geometric,calc}
\usepackage{tabularx,booktabs,array}
\usepackage{enumitem}
\usepackage{parskip}
\usepackage{fancyhdr}
\usepackage{setspace}
\usepackage{xcolor}
\usepackage[hidelinks]{hyperref}
\definecolor{accent}{HTML}{008080}
\definecolor{darkaccent}{HTML}{004D4D}
\definecolor{textgray}{HTML}{333333}
\definecolor{muted}{HTML}{666666}
\definecolor{thinrule}{HTML}{ABD0CE}
\color{textgray}
\onehalfspacing
\setlength{\parindent}{0pt}
\setlength{\parskip}{5pt}
\setlist[itemize]{leftmargin=5mm,itemsep=1pt,topsep=3pt}
\setlist[enumerate]{leftmargin=7mm,itemsep=5pt,topsep=3pt}
\pagestyle{fancy}
\fancyhf{}
\fancyhead[L]{\footnotesize\color{muted} Логарифмические неравенства}
\fancyhead[R]{\footnotesize\color{muted} Тимофей Васильченко}
\fancyfoot[C]{\footnotesize\thepage}
\renewcommand{\headrulewidth}{0.35pt}
\renewcommand{\headrule}{\hbox to\headwidth{\color{accent}\leaders\hrule height \headrulewidth\hfill}}
\newcommand{\topic}[1]{\vspace{0.65em}{\large\color{darkaccent}\bfseries #1}\par\vspace{0.15em}{\color{thinrule}\rule{\linewidth}{0.5pt}}\vspace{0.35em}}
\newcommand{\subtopic}[1]{\vspace{0.35em}{\color{darkaccent}\bfseries #1}\par}
\newcommand{\idea}[1]{{\color{darkaccent}\textbf{Ключевая идея.}} #1}
\newcommand{\example}[1]{{\color{darkaccent}\textbf{Пример #1.}}}
\begin{document}

% ---------------------------------------------------------------------
% Титульный лист
% ---------------------------------------------------------------------
\begin{titlepage}
\thispagestyle{empty}
\centering
\vspace*{24mm}
{\normalsize\color{darkaccent} МАТЕМАТИКА\par}
\vspace{18mm}
{\LARGE\color{textgray} Чек-лист\par}
\vspace{12mm}
{\Large\color{accent} Логарифмические неравенства:\par}
{\Large\color{accent} ОДЗ, замена и модуль\par}
\vspace{17mm}
{\normalsize Подготовка к ЕГЭ по профильной математике\par}
\vfill
{\large Тимофей Васильченко\par}
\vspace{10mm}
{\normalsize 31 октября 2026 г.\par}
\vspace{19mm}
{\small Лёвин Артём Александрович\par}
{\small эксклюзивно для Тимофея Васильченко\par}
\vspace{15mm}
\begin{tikzpicture}[x=1cm,y=1cm]
  \draw[accent,line width=0.85pt] (0,0) .. controls (2,0.78) and (3.9,-0.50) .. (6.0,0)
                                    .. controls (8.0,0.68) and (9.3,-0.48) .. (11.2,0.04);
\end{tikzpicture}
\vspace{7mm}
\end{titlepage}
\setcounter{page}{1}

\topic{1. С чего начинается логарифмическое неравенство}
\idea{До преобразований запишите область допустимых значений (ОДЗ).}

Логарифм $\log_a b$ определён в действительных числах тогда и только тогда, когда
\[
\boxed{a>0,\qquad a\ne1,\qquad b>0.}
\]
Если основание или аргумент содержит $x$, все три ограничения проверяются для выражений с переменной.

\example{1} Найдём ОДЗ выражения $\log_{x+1}\bigl((x-2)^2\bigr)$.
\[
\begin{cases}
x+1>0,\\[-1mm]
x+1\ne1,\\[-1mm]
(x-2)^2>0
\end{cases}
\quad\Longleftrightarrow\quad
\begin{cases}
x>-1,\\[-1mm]
x\ne0,\\[-1mm]
x\ne2.
\end{cases}
\]
У чётной степени результат неотрицательный; нуль нужно исключать, поскольку аргумент логарифма строго положителен.

\topic{2. Как избавляться от степеней под логарифмом}
При $a>0$, $a\ne1$ и $B\ne0$ имеем
\[
\boxed{\log_a(B^{2n})=2n\log_a|B|}\qquad(n\in\mathbb N).
\]
Модуль необходим: исходное $B$ может быть отрицательным, а $B^{2n}$ всё равно положительно.

\example{2} При $x\ne4$:
\[
\log_8\bigl((4-x)^4\bigr)
=4\log_8|4-x|
=\frac43\log_2|4-x|,
\]
\[
\log_{0{,}5}\bigl((4-x)^8\bigr)
=\log_{2^{-1}}\bigl((4-x)^8\bigr)
=-8\log_2|4-x|.
\]
Здесь использовано правило замены основания:
\[
\log_{a^k}B=\frac1k\log_a B\qquad(k\ne0).
\]
\subtopic{Проверьте себя}
Чётная степень требует модуля, основание $\tfrac12$ меняет знак при переходе к основанию $2$, а найденная в начале ОДЗ сохраняется до ответа.

\clearpage
\topic{3. Главный приём: заменить повторяющийся логарифм}
\idea{Сначала приведите степени и основания к одному виду. После этого вводите замену $t$.}

\example{3} Решим неравенство
\[
9\bigl(\log_8((4-x)^4)\bigr)^2
+5\log_{0{,}5}((4-x)^8)\le56.
\]
\textbf{Шаг 1. ОДЗ.} Аргументы логарифмов строго положительны при $x\ne4$.

\textbf{Шаг 2. Приведение к основанию $2$.}
\[
9\left(\frac43\log_2|4-x|\right)^2
-40\log_2|4-x|\le56.
\]
Важно: при возведении в квадрат возводится и коэффициент $\tfrac43$.
\[
16\bigl(\log_2|4-x|\bigr)^2
-40\log_2|4-x|-56\le0.
\]

\textbf{Шаг 3. Замена.} Положим $t=\log_2|4-x|$. Разделим на $8>0$:
\[
2t^2-5t-7\le0.
\]
Дискриминант $D=(-5)^2-4\cdot2\cdot(-7)=81$; корни:
\[
t_1=\frac{5-9}{4}=-1,
\qquad t_2=\frac{5+9}{4}=\frac72.
\]
Коэффициент при $t^2$ положительный, поэтому
\[
\boxed{-1\le t\le\frac72.}
\]
\begin{center}
\begin{tikzpicture}[>=Stealth,scale=0.94]
 \draw[->,line width=.7pt] (0,0)--(11.6,0) node[right] {$t$};
 \draw[accent,line width=2.6pt] (3.2,0)--(8.9,0);
 \fill[accent] (3.2,0) circle (2.2pt) (8.9,0) circle (2.2pt);
 \node[below=5pt] at (3.2,0) {$-1$};
 \node[below=5pt] at (8.9,0) {$\tfrac72$};
 \node[above=6pt,text=muted] at (6.05,0) {$2t^2-5t-7\le0$};
\end{tikzpicture}
\end{center}
\textbf{Шаг 4. Обратная замена.} Так как основание $2>1$, направление знаков сохраняется:
\[
-1\le\log_2|4-x|\le\frac72
\quad\Longleftrightarrow\quad
\frac12\le|4-x|\le8\sqrt2.
\]

\clearpage
\topic{4. Как раскрывать модуль и собирать ответ}
Для $c\ge0$ действуют два перехода:
\[
\boxed{|F(x)|\ge c\ \Longleftrightarrow\
\begin{cases}
F(x)\ge c\quad\text{или}\\ F(x)\le-c;
\end{cases}}
\]
\[
\boxed{|F(x)|\le c\ \Longleftrightarrow\
\begin{cases}
F(x)\ge-c,\\ F(x)\le c.
\end{cases}}
\]
В первом случае возникает \textbf{совокупность} (объединение решений), во втором --- \textbf{система} (пересечение решений). Для строгих неравенств границы становятся строгими.

\subtopic{Завершим пример 3}
Из $\tfrac12\le|4-x|\le8\sqrt2$ одновременно требуются два условия:
\[
\underbrace{|4-x|\ge\frac12}_{\text{совокупность}}
\quad\text{и}\quad
\underbrace{|4-x|\le8\sqrt2}_{\text{система}}.
\]
Первое даёт $x\le\tfrac72$ или $x\ge\tfrac92$.
Второе даёт $4-8\sqrt2\le x\le4+8\sqrt2$.
Пересекаем эти множества:
\[
\boxed{x\in[4-8\sqrt2;\,\tfrac72]
\ \cup\ [\tfrac92;\,4+8\sqrt2].}
\]
Точка $x=4$ в ответ не входит автоматически: она лежит между полученными промежутками.

\topic{5. Логарифмы с переменным основанием}
Если основания одинаковы, но зависят от $x$, сначала найдите ОДЗ:
\[
a(x)>0,\qquad a(x)\ne1,\qquad F(x)>0,\qquad G(x)>0.
\]
Затем можно использовать равносильный переход:
\[
\boxed{\log_{a(x)}F(x)\ge\log_{a(x)}G(x)
\ \Longleftrightarrow\
(a(x)-1)\bigl(F(x)-G(x)\bigr)\ge0}
\]
\textit{при указанных выше ограничениях.} Множитель $a(x)-1$ учитывает, возрастает логарифм по аргументу или убывает.

\clearpage
\topic{6. Переменное основание: пример с единицей}
\example{4} Решим неравенство
\[
\log_{x+1}(x^2+x-6)\ge1.
\]
\textbf{Шаг 1. ОДЗ.}
\[
\begin{cases}
x+1>0,\\[-1mm]
x+1\ne1,\\[-1mm]
x^2+x-6>0.
\end{cases}
\]
Разложим аргумент: $x^2+x-6=(x+3)(x-2)$. Совместно с $x>-1$ получим $x>2$. В частности, основание $x+1>1$.

\textbf{Шаг 2. Представим единицу логарифмом с тем же основанием:}
\[
1=\log_{x+1}(x+1).
\]
Поскольку на ОДЗ основание больше $1$, логарифмическая функция возрастает:
\[
x^2+x-6\ge x+1
\quad\Longleftrightarrow\quad x^2-7\ge0.
\]
\textbf{Шаг 3. Пересечём решение с ОДЗ:}
\[
\left(x\le-\sqrt7\quad\text{или}\quad x\ge\sqrt7\right)
\quad\text{и}\quad x>2.
\]
\[
\boxed{x\in[\sqrt7;+\infty).}
\]
Тот же вывод даёт общий приём $(a-1)(F-G)\ge0$:
\[
\begin{cases}
x>2,\\[-1mm]
x(x^2-7)\ge0.
\end{cases}
\]

\topic{7. Типичные ошибки и короткая памятка}
Не теряйте модуль при вынесении чётной степени; учитывайте отрицательный знак при переходе от основания $\tfrac12$ к $2$; отличайте пересечение решений системы от объединения решений совокупности. Для переменного основания проверяйте его положение относительно $1$. Всегда сверяйте ответ с ОДЗ.

\clearpage
\topic{Алгоритм: решение логарифмического неравенства}
\vspace{5mm}
\begin{center}
\tikzset{
 start/.style={ellipse,draw=darkaccent,very thick,align=center,minimum width=4.4cm,minimum height=1.1cm,fill=white},
 act/.style={rectangle,rounded corners=2pt,draw=thinrule,thick,align=center,text width=9.7cm,minimum height=1.18cm,inner sep=5pt,fill=white},
 choice/.style={diamond,aspect=2.1,draw=accent,thick,align=center,text width=4.0cm,inner sep=2.5pt,fill=white},
 arr/.style={-{Stealth[length=2.3mm]},line width=.8pt,draw=accent}
}
\begin{tikzpicture}[every node/.style={font=\small},>=Stealth]
\node[start] (s) at (0,0) {Начало};
\node[act] (odz) at (0,-1.85) {Запишите ОДЗ всех логарифмов};
\node[act] (simp) at (0,-3.70) {Приведите степени и основания к удобному виду; при чётной степени сохраните модуль};
\node[choice] (which) at (0,-6.00) {Есть повторяющийся логарифм?};
\node[act,text width=4.25cm,minimum height=1.65cm] (sub) at (-3.0,-8.8) {Введите $t$ и решите неравенство относительно $t$};
\node[act,text width=4.25cm,minimum height=1.65cm] (base) at (3.0,-8.8) {Сравните логарифмы с учётом основания};
\node[act] (back) at (0,-11.20) {Вернитесь к $x$; раскройте модуль, найдите пересечения и объединения};
\node[act] (check) at (0,-13.10) {Сопоставьте полученные интервалы с ОДЗ; проверьте границы};
\node[start] (end) at (0,-15.00) {Ответ};
\draw[arr] (s)--(odz);
\draw[arr] (odz)--(simp);
\draw[arr] (simp)--(which);
\draw[arr] (which.south west) -- node[above left,sloped,font=\footnotesize]{да} (sub.north);
\draw[arr] (which.south east) -- node[above right,sloped,font=\footnotesize]{нет} (base.north);
\draw[arr] (sub.south) -- ([xshift=-3cm]back.north);
\draw[arr] (base.south) -- ([xshift=3cm]back.north);
\draw[arr] (back)--(check);
\draw[arr] (check)--(end);
\end{tikzpicture}
\end{center}
\clearpage

\topic{Тренировочный блок --- Путь А}
Выполняйте задания по порядку. Учитывайте ОДЗ, если выражение содержит логарифм с переменной, и сохраняйте точные значения корней.

\subtopic{Базовый уровень}
\begin{enumerate}
\item Найдите область допустимых значений выражения
\[\log_{x+2}\bigl((x-1)^2\bigr).\]
\item При $x\ne3$ выразите через $\log_2|x-3|$:
\[\log_8\bigl((x-3)^4\bigr).\]
\item При $x\ne-1$ выразите через $\log_2|x+1|$:
\[\log_{0{,}5}\bigl((x+1)^6\bigr).\]
\item Решите неравенство:
\[2t^2-5t-3\le0.\]
\end{enumerate}
\subtopic{Средний уровень}
\begin{enumerate}[resume]
\item Решите неравенство $|x-4|\ge\dfrac32$.
\item Решите двойное неравенство $1\le|x-3|\le4$.
\item Решите логарифмическое неравенство
\[\log_2|x-1|\le2.\]
\end{enumerate}
\clearpage
\topic{Тренировочный блок --- продолжение}
\subtopic{Уровень выше среднего}
\begin{enumerate}[resume]
\item Решите неравенство (сначала сведите его к квадратному заменой):
\[
9\bigl(\log_8((x-2)^4)\bigr)^2
+3\log_{0{,}5}((x-2)^8)\le16.
\]
\item Решите неравенство с переменным основанием:
\[
\log_{x+1}(x+4)\ge\log_{x+1}(2x+1).
\]
\end{enumerate}
\subtopic{Сложный уровень}
\begin{enumerate}[resume]
\item Решите неравенство:
\[
\log_{x+1}(x^2+x-6)\ge1.
\]
\end{enumerate}
\vfill
{\footnotesize\color{muted} Подсказка к оформлению: в ответах на неравенства используйте объединения промежутков и квадратные скобки только там, где граница действительно входит в решение.}
\clearpage

\topic{Ответы к тренировочному блоку}
\renewcommand{\arraystretch}{1.65}
\begin{tabularx}{\linewidth}{@{}>{\centering\arraybackslash}p{13mm} X@{}}
\toprule
\textbf{№} & \textbf{Ответ}\\
\midrule
1 & $x>-2,\quad x\ne-1,\quad x\ne1$.\\
2 & $\dfrac43\log_2|x-3|$ при $x\ne3$.\\
3 & $-6\log_2|x+1|$ при $x\ne-1$.\\
4 & $t\in\left[-\dfrac12;\,3\right]$.\\
5 & $x\in\left(-\infty;\,\dfrac52\right]\cup\left[\dfrac{11}2;\,+\infty\right)$.\\
6 & $x\in[-1;2]\cup[4;7]$.\\
7 & $x\in[-3;1)\cup(1;5]$.\\
8 & $x\in\left[-2;\,2-\dfrac{1}{\sqrt2}\right]\cup\left[2+\dfrac{1}{\sqrt2};\,6\right]$.\\
9 & $x\in(0;3]$.\\
10 & $x\in[\sqrt7;+\infty)$.\\
\bottomrule
\end{tabularx}

\vspace{12mm}
\topic{Итоговая памятка}
\begin{tabularx}{\linewidth}{@{}>{\raggedright\arraybackslash}p{54mm} X@{}}
\toprule
\textbf{Признак в условии} & \textbf{Действие}\\
\midrule
Логарифм с переменной & Выписать ОДЗ для основания и аргумента.\\
Чётная степень под логарифмом & Вынести показатель, поставив модуль.\\
Основание $8$ или $0{,}5$ & Привести к основанию $2$.\\
Повторяется один логарифм & Ввести замену $t$.\\
$|F|\ge c$ или $|F|\le c$ & Различать совокупность и систему.\\
Основание зависит от $x$ & Учитывать знак $a(x)-1$.\\
\bottomrule
\end{tabularx}
\end{document}
